\begin{korollar}{Norm und duale Norm in Hilbertr�umen}
                {NormUndDualeNormImHilbertraum}
Seien $(H, \sprod\cdot{\cdot\cdot})$ ein $\K$--Hilbertraum, $x \in H$ und
$K \in \meng{ \kran H 0 1,\kabg H 0 1}$.\\
Dann gilt $\norm x_H = \sup\menge{\abs{\sprod x y}}{y \in K}$.\\
Im Falle $\K = \R$ gilt au�erdem $\norm x_H = \sup\menge{\sprod x y}{y \in K}$.
\end{korollar}
\begin{proof}
Nach dem aus der Funktionalanalysis bekannten Darstellungssatz von
Fr�chet--Riesz ist $\iota : H \rightarrow H', x \mapsto \sprod\cdot x$
eine isometrische, antilineare Bijektion.\\
Daher ist $\menge{\abs{\sprod x y}}{y \in K} 
= \menge{\abs{\vi(x)}}{\vi \in \iota(K)}$.\\
Da $\iota$ eine bijektive Isometrie ist, gilt 
$\iota(K) \in \meng{\kabg{H'}01, \kran{H'}01}$.\\
Somit folgt der Rest der Behauptung direkt aus \ref{NormUndDualeNorm}.
\end{proof}